\chapter{The Impact of DHT Routing Geometry on Resilience and Proximity}
\paperinfo{K. Gummadi, R. Gummadi, S. Gribble, S. Ratnasamy, S. Shenker, I. Stoica (U. Washington, USC, Intel Research, ICSI, UC Berkeley)}{ACM SIGCOMM 2003}

\begin{keybox}[The paper in three sentences]
Many DHTs (Chord, CAN, Pastry, Kademlia, Viceroy, \dots) all give $O(\log n)$ hops with $O(\log n)$ neighbours, so to compare them the paper looks
at their underlying \key{routing geometry} (tree, hypercube, butterfly, ring, XOR, hybrid) and how much \key{flexibility} each geometry
leaves in choosing \key{neighbours} and \key{next hops}. Flexibility decides how well a DHT routes when nodes fail (\key{static
resilience}) and how well it can use short-latency links (\key{proximity}). The \key{ring} turns out to be the most flexible, so it is the most
resilient and as good as any for latency: ``why not use rings?''
\end{keybox}

% ------------------------------------------------------------------
\section{The question: how should we compare DHTs?}

A DHT gives a hash-table interface over many nodes: \texttt{lookup(key)} finds the node responsible. Dozens of routing algorithms had been
proposed, each in its own paper, each claiming to be best. Two ways to compare them:
\begin{itemize}
  \item \textbf{Bake-off}: run each complete system and measure. This ranks today's systems but teaches little about how to design a better one.
  \item \textbf{Component analysis} (this paper): split each design into separate choices and ask which choice causes which property.
\end{itemize}

\textbf{Two kinds of choices}:
\begin{itemize}
  \item \key{Routing level}: \key{neighbour selection} (which nodes go in my routing table) and \key{route selection} (which neighbour to forward to).
  \item \key{System level}: everything else (caching, replication, iterative vs recursive lookups); mostly independent of routing.
\end{itemize}

\textbf{Algorithm vs geometry}. The \key{algorithm} is the exact rule for neighbours and next hops; tiny changes create a new algorithm. The
\key{geometry} is the shape behind it (ring, tree, \dots); it constrains what algorithms are possible. Small algorithm tweaks do not change the
geometry. The paper studies geometry, because there are few geometries but endless algorithms.

\textbf{Why state and hop count do not separate them}: all geometries here give $O(\log n)$ hops with $O(\log n)$ neighbours (the butterfly even
manages $O(1)$ neighbours). The difference lies in \key{flexibility}: how much freedom the geometry leaves, once fixed, in picking neighbours and routes.

\begin{keybox}[Two kinds of flexibility]
\begin{itemize}
  \item \key{Neighbour selection flexibility}: for each routing-table slot, how many nodes are allowed? More choice $\Rightarrow$ you can pick the
        \emph{nearest} (lowest latency) one.
  \item \key{Route selection flexibility}: for a given destination, how many of my neighbours make progress? More choice $\Rightarrow$ if one is
        dead you can use another, or pick the nearest.
\end{itemize}
Also important: \key{sequential neighbours}, neighbours that make progress towards \emph{every} destination (Chord's successors, Pastry's leaf set).
They are the last resort when everything else fails.
\end{keybox}

% ------------------------------------------------------------------
\section{The six geometries, with a tiny example}

Assume $n$ nodes with $\log n$-bit IDs, the ID space fully used. Example: $n = 8$, 3-bit IDs \texttt{000}--\texttt{111}; route from
\texttt{000} to \texttt{111}.

\subsection{Tree (PRR, Tapestry; Pastry's main mode)}
\textbf{Picture}: IDs are the leaves of a binary tree. Distance between two nodes $=$ height of their smallest common subtree (how many trailing
bits differ, counted from the first differing bit).
\textbf{Neighbours}: one neighbour at each distance $i$: a node that shares my prefix up to that level but differs in the next bit.
\textbf{Routing}: fix the \emph{highest} differing bit first, so the matching prefix grows by one bit per hop.

\begin{center}
\small
\begin{tabular}{@{}lll@{}}
\toprule
Slot for node \texttt{000} & Allowed nodes & Choices\\
\midrule
distance 1 (differ in last bit) & \texttt{001} & 1\\
distance 2 & \texttt{010}, \texttt{011} & 2\\
distance 3 (differ in first bit) & \texttt{100}, \texttt{101}, \texttt{110}, \texttt{111} & 4\\
\bottomrule
\end{tabular}
\end{center}
\begin{itemize}
  \item \textbf{Neighbour flexibility: large}. Slot $i$ has $2^{i-1}$ choices: the low-order bits are free. In total about $n^{(\log n)/2}$ possible tables.
  \item \textbf{Route flexibility: none}. To reach \texttt{111}, the first bit must be fixed first, and exactly one neighbour (the distance-3 one) does that.
        If it is dead, routing fails.
\end{itemize}

\subsection{Hypercube (CAN with $\log n$ dimensions)}
\textbf{Neighbours}: the $\log n$ nodes that differ from me in exactly one bit. For \texttt{000}: \texttt{100}, \texttt{010}, \texttt{001}.
\textbf{Routing}: fix any differing bit, \emph{in any order}. A neighbour differs from me in one bit only, so fixing one bit never breaks bits already fixed.
\begin{itemize}
  \item \textbf{Neighbour flexibility: none}. Each slot has exactly one allowed node.
  \item \textbf{Route flexibility: large}. From \texttt{000} to \texttt{111}: 3 choices for the first hop, 2 for the second, 1 for the last, so $3! = 6$
        shortest paths; in general about $(\log n)!$, all of the same length.
\end{itemize}
The tree and the hypercube are opposites: one has only neighbour freedom, the other only route freedom.

\subsection{Butterfly (Viceroy)}
Nodes are arranged in $\log n$ stages; stage $i$ fixes bit $i$. Viceroy gets $O(\log n)$ hops with only \key{constant} state. But routing goes up to stage 1,
down the stages, and then uses successor/predecessor links for the last stretch. \textbf{No flexibility} in neighbours or routes. The authors suspect this is the
price of constant state.

\subsection{Ring (Chord)}
IDs on a circle; distance $=$ clockwise distance. Chord's rule says finger $i$ is \emph{exactly} the node closest to $a + 2^i$. But nothing in the ring
geometry needs that exactness: routing still takes $O(\log n)$ hops if finger $i$ is \key{any node in $[a + 2^i, a + 2^{i+1})$}.
\begin{itemize}
  \item \textbf{Neighbour flexibility: large}: $2^i$ choices for slot $i$, like the tree.
  \item \textbf{Route flexibility: large}: the path to a destination at distance $d$ is made of spans $2^i$ for the 1-bits of $d$, and \emph{these spans can be taken in
        any order}. From 0 to 7 ($d = 111_2$): spans 4, 2, 1 in any order, so 6 paths, all of 3 hops. On top of that, a hop of span $2^i$ can be replaced by two hops
        of span $2^{i-1}$, giving many \emph{longer} backup paths too.
  \item \textbf{Sequential neighbours: natural} (successors), since there is one global order around the circle.
\end{itemize}
So the ring has \emph{both} kinds of flexibility.

\subsection{XOR (Kademlia)}
Distance between two IDs $=$ their bitwise XOR, read as a number. Neighbour $i$ is any node at XOR distance in $[2^i, 2^{i+1})$: exactly the same allowed tables as the
tree, and with no failures greedy routing follows the same paths as the tree (fixing bits from the left). Under failures it is better than the tree: if the
highest bit cannot be fixed, it can fix a lower bit and still reduce the XOR distance. But a lower bit fixed early may be undone later, so these alternative
paths are \emph{longer}.

\subsection{Hybrid (Pastry)}
IDs are both tree leaves and points on a ring. Pastry routes on the tree by default and falls back to the ring (its leaf set) when tree routing fails. Neighbour
flexibility as the tree; extra routes come only from the fallback and are not always $O(\log n)$ long.

\subsection{Summary table}
\begin{center}
\small
\begin{tabular}{@{}lcccccc@{}}
\toprule
Property & Tree & Hypercube & Ring & Butterfly & XOR & Hybrid\\
\midrule
Neighbour selection (tables) & $n^{\log n/2}$ & 1 & $n^{\log n/2}$ & 1 & $n^{\log n/2}$ & $n^{\log n/2}$\\
Route selection, optimal paths & 1 & $c_1(\log n)!$ & $c_1(\log n)!$ & 1 & 1 & 1\\
Route selection, longer paths & -- & -- & $2c_2(\log n)!$ & -- & $c_2(\log n)!$ & $c_2(\log n)!$\\
Natural sequential neighbours & no & no & \key{yes} & no & no & fallback only\\
\bottomrule
\end{tabular}
\end{center}
\textbf{Why the ring and hypercube have twice the route flexibility of XOR and hybrid}: on the ring, for a destination at distance in $[2^i, 2^{i+1})$, \emph{all} of my
first $i$ neighbours make progress. In XOR/hybrid, my $i$-th neighbour makes progress only if its $i$-th bit matches the destination's, which is true half the time.

% ------------------------------------------------------------------
\section{Static resilience}

Resilience to failures has three parts:
\begin{enumerate}
  \item \textbf{Data replication}: keep copies, so data is not lost (assumed adequate here).
  \item \textbf{Routing recovery}: algorithms that refill routing tables with live nodes (take time).
  \item \key{Static resilience}: \emph{before} recovery has done anything, with dead entries simply removed, how much routing still works?
\end{enumerate}
Static resilience tells you how fast recovery has to be: a DHT with low static resilience needs very fast recovery.

\textbf{Experiment}: 65,536 nodes, the same table size for every geometry (except the butterfly, whose state is fixed). With no failures the hop counts are similar (about
7.7 average for all; butterfly 21.4). Then kill a random fraction of nodes, delete them from tables, try to route between every pair of live nodes, and measure
\% failed paths and \% increase in path length.

\paperfig[0.92]{dg_fig1}{Left: \% of paths that fail as nodes fail. Right: \% increase in hop count of the paths that succeed. (No sequential neighbours.)\label{fig:dgres}}

\textbf{Results at 30\% failed nodes}:
\begin{center}
\small
\begin{tabular}{@{}lll@{}}
\toprule
Geometry & Failed paths & Why\\
\midrule
Tree, Butterfly & about \key{90\%} & no route flexibility: one dead hop kills the path\\
XOR, Hybrid & about \key{20\%} & about half the route flexibility of ring/hypercube\\
Ring, Hypercube & \key{under 7\%} & the most route flexibility\\
\bottomrule
\end{tabular}
\end{center}
\begin{itemize}
  \item \textbf{Path stretch}: the hypercube barely stretches (all its alternatives are equally short); the ring stretches somewhat (it also uses longer alternatives);
        XOR and hybrid stretch a lot (their alternatives are only the longer kind). The ring is slightly more resilient than the hypercube \emph{because} it also has
        those longer paths.
  \item \textbf{Adding 16 sequential neighbours} to every geometry: \key{no path fails} even at 30\% failures. The ring (with natural sequential neighbours) does best, but
        paths get longer.
  \item \textbf{Are sequential neighbours better than just more ordinary neighbours?} For the ring, at the same table size, sequential neighbours give a bit more success at
        high failure rates, but ordinary neighbours give much shorter paths.
\end{itemize}

\begin{keybox}[Resilience conclusion]
Static resilience follows \key{route selection flexibility}. The ring has the most, so it is the most resilient; tree and butterfly have none and collapse.
\end{keybox}

% ------------------------------------------------------------------
\section{Proximity: path latency}

Hop count is not latency. Each overlay hop may cross an ocean, so an ``efficient'' 7-hop route can be 10 times slower than the direct IP path. Ways to
use proximity:
\begin{itemize}
  \item \key{PNS, Proximity Neighbour Selection}: when filling a routing-table slot, among the allowed nodes pick the \emph{closest} (lowest latency). Needs
        \emph{neighbour} flexibility (tree, ring, XOR; not hypercube).
  \item \key{PRS, Proximity Route Selection}: with a fixed table, when forwarding, among neighbours that make progress pick the closest. Needs \emph{route}
        flexibility (ring, hypercube, XOR; not tree).
  \item PIS (choose IDs by location): not studied, since it breaks load balance and correlates failures.
\end{itemize}
Real PNS cannot check all $2^i$ candidates of a slot, so \key{PNS($K$)} samples $K$ of them and picks the closest.

\textbf{Modelling the Internet's latencies.} Results depend heavily on latencies, and topology generators like GT-ITM assign them arbitrarily. The authors conjecture
that the \key{latency distribution seen from one typical node} is enough: treat each candidate's latency as a random draw from it. They check this against simulation and
it matches well. Real measured distributions differ from GT-ITM: in reality at least 10\% of paths are more than twice the median latency; in GT-ITM none are. So results
from GT-ITM can mislead.

\subsection{Which is better, PNS or PRS?}
For XOR and ring (which allow both): PRS helps somewhat, \key{PNS helps much more}, and PNS+PRS adds only a little over PNS.

\textbf{Why PNS wins (counting choices).} Route to a destination at distance between $2^i$ and $2^{i+1}$ on a ring:
\begin{center}
\small
\begin{tabular}{@{}lll@{}}
\toprule
 & What you choose from & Number of options\\
\midrule
PNS & the $i$-th neighbour can be \emph{any} of the $2^i$ nodes in its range: pick the nearest & $2^i$ (e.g.\ 1024 for $i = 10$)\\
PRS & any of my first $i$ neighbours, each fixed in advance & $i$ (e.g.\ 10)\\
\bottomrule
\end{tabular}
\end{center}
The minimum latency over 1024 random draws is much lower than over 10. And PNS is paid once (at table-building time), not per lookup.

\subsection{Does the geometry matter beyond allowing PNS/PRS?}
No. PNS Ring $\approx$ PNS Tree; PRS Ring $\approx$ PRS Hypercube; PNS+PRS Ring $\approx$ PNS+PRS XOR. What matters is \emph{whether} the geometry permits PNS.
So the hypercube, which cannot do PNS, is worse for latency.

\textbf{How good is it in absolute terms?} Median latency (ms) with the Virginia latency distribution (16,384 nodes):
\begin{center}
\begin{tabular}{@{}lcc@{}}
\toprule
Scheme & Hops & Median latency\\
\midrule
Direct Internet path & 1 & 102\\
Plain XOR & 7 & 1036 (about 10$\times$)\\
PRS XOR & 8 & 770\\
PNS XOR & 7 & 139\\
PNS+PRS XOR & 7 & \key{136 (under 2$\times$)}\\
\bottomrule
\end{tabular}
\end{center}

% ------------------------------------------------------------------
\section{Local convergence}

\key{Local convergence}: two messages from \emph{nearby} sources to the same key should meet at a node near the sources, before travelling far. It helps overlay
multicast (efficient trees), caching along lookup paths (a nearby node's cached copy is found), and server selection.

\textbf{Experiment}: an isolated domain of $m$ nodes, all close to each other and far from the other $n - m$. Every node in the domain looks up the same outside key.
Count \key{exit points}: domain nodes that forward the message out of the domain. Ideal $= 1$.
\begin{itemize}
  \item \textbf{PNS and PNS+PRS} are almost optimal. \textbf{PRS alone} barely helps: with 16 neighbours out of 65,536 nodes, it is unlikely any neighbour is inside the domain,
        so there is nothing nearby to choose.
  \item With realistic \textbf{PNS(16)} sampling, PNS alone is about as weak as PRS, and \textbf{PNS(16)+PRS} is clearly better. So PRS matters when PNS can only be approximated.
  \item PNS Tree, PNS+PRS Ring and PNS+PRS XOR perform about the same: again, the geometry matters only through what it permits.
\end{itemize}

% ------------------------------------------------------------------
\section{Conclusions}
\begin{enumerate}
  \item \textbf{Analyse components, not whole systems.}
  \item \textbf{Geometry is fundamental}: it constrains everything else, and there are few geometries, so agreement on geometry can come first.
  \item \textbf{Flexibility is what differs} between geometries (apart from the butterfly's constant state), and it drives resilience, latency and convergence.
  \item \textbf{Ring and XOR} are the only geometries that allow both PNS and PRS.
  \item \textbf{Why not the ring?} Most flexible, natural sequential neighbours, best static resilience, as good as any on latency and convergence. The authors had
        expected a more complex geometry to win; it did not.
\end{enumerate}

\begin{keybox}[Exam one-liner]
All DHT geometries give $O(\log n)$ hops with $O(\log n)$ state, so the paper compares their \emph{flexibility}: freedom in choosing neighbours (tree, ring, XOR: $2^{i}$
options per slot; hypercube, butterfly: none) and in choosing next hops (hypercube, ring: about $(\log n)!$ routes; tree, butterfly: one). Static resilience follows route
flexibility (at 30\% failures: tree/butterfly about 90\% of paths fail, ring/hypercube under 7\%; sequential neighbours fix nearly all failures). For latency, Proximity
Neighbour Selection beats Proximity Route Selection ($2^i$ vs $i$ choices), bringing paths to under 2$\times$ the IP latency, and geometry matters only through whether it
allows PNS. The ring has both flexibilities plus natural sequential neighbours, so ``why not the ring?''
\end{keybox}

\section{Revision checklist}
\begin{itemize}
  \item Routing level (neighbour/route selection) vs system level; algorithm vs geometry.
  \item Flexibility definitions; sequential neighbours.
  \item Tree: neighbour freedom $2^{i-1}$, route freedom 1. Hypercube: neighbour 1, routes $(\log n)!$ (any bit order). Butterfly: none, constant state. Ring: both, spans in
        any order + longer paths, natural successors. XOR: tree-like table, alternatives longer. Hybrid: tree + ring fallback.
  \item Why ring/hypercube have 2$\times$ route flexibility of XOR/hybrid.
  \item Static resilience definition; results at 30\%; stretch; 16 sequential neighbours $\to$ no failures.
  \item PNS vs PRS vs PIS; PNS($K$); latency-distribution conjecture; GT-ITM caveat.
  \item PNS $\gg$ PRS ($2^i$ vs $i$); geometry matters only via PNS/PRS support; XOR plain 1036\,ms vs PNS 139\,ms vs IP 102\,ms.
  \item Local convergence: exit points; PNS near optimal; PRS helps when PNS is sampled (PNS(16)).
\end{itemize}
